\documentclass[11pt]{article}
\usepackage{amsmath,amssymb,booktabs}
\begin{document}

\section*{Step 219: Analytical Pattern for \(\Phi(\sigma,\ell)\)}

\subsection*{Definition}

For the large-\(T\) Gaussian wavepacket sequence, define
\[
\Phi(\sigma,\ell)=\lim_{T\to\infty}\|C_\ell u_T\|,
\qquad C_\ell=(I-P_\infty)M_{m_\ell}P_\infty.
\]
Numerically this is evaluated at \(T=10000\).

\subsection*{\(\ell\)-Sweep at \(\sigma=1\)}

The values rise from \(0.1309\) at \(\ell=0.1\) to a plateau near \(0.279\) for \(\ell\in[1.5,2]\), then decay for larger shifts:
\[
\Phi(1,\log 2)=0.2638458688,\qquad
\Phi(1,2)=0.2778803484,\qquad
\Phi(1,5)=0.0011147046.
\]

\subsection*{\(\sigma\)-Sweep at \(\ell=\log 2\)}

\[
\Phi(0.25,\log 2)=0.2931,\quad
\Phi(0.5,\log 2)=0.3512,\quad
\Phi(1,\log 2)=0.2638,\quad
\Phi(2,\log 2)=0.0470,\quad
\Phi(5,\log 2)=0.0012.
\]
Thus the dependence on \(\sigma\) is non-monotone.

\subsection*{Band-Shift Model}

Ignoring finite PSWF correction, \(P_\infty\) is modeled by a sinc cutoff, i.e. a spectral projection onto a band.  The modulation \(e^{i\ell\tau}\) shifts the Gaussian spectrum by \(\ell\).  The amount moved from the high-pass region into the low-pass region is modeled by
\[
\Phi(\sigma,\ell)^2\approx
\frac12\left(\operatorname{erf}(\sigma b(\ell))-\operatorname{erf}(\sigma a(\ell))\right),
\]
where
\[
a(\ell)=
\begin{cases}
1,& 0\le \ell\le 2,\\
\ell-1,& \ell>2,
\end{cases}
\qquad
b(\ell)=\ell+1.
\]
For the \(\ell\)-sweep at \(\sigma=1\), this model has RMSE \(2.40\cdot10^{-3}\).  With fitted scale and effective width, the RMSE improves to \(1.47\cdot10^{-3}\).  For the \(\sigma\)-sweep at \(\ell=\log 2\), the same model has RMSE \(1.47\cdot10^{-3}\).

\subsection*{Verdict}

The data show a clean band-shift error-function pattern, but not a closed form at the \(10^{-4}\) residual standard:
\[
\boxed{\mathrm{V\_essential\_norm\_ell\_dependence\_numerical\_pattern}}.
\]

\end{document}
